\section{Conclusion}
\label{sec:conclusion}

We began by observing that adapter selection in multi-task models typically requires validation data that is unavailable in zero-shot settings. This work demonstrates that the validation data requirement can be eliminated---instruction prefixes alone contain sufficient signal for effective routing.

\subsection{Summary}

We introduced Instruction-Prefix-Conditioned Routing (IPCR), a minimal approach for zero-shot adapter selection using frozen MiniLM embeddings and a linear probe. Our experiments on FLAN task families establish three results:

\begin{enumerate}
    \item \textbf{Intrinsic alignment exists:} Instruction embeddings are linearly separable by task family with macro-F1 = 0.995, demonstrating that instruction semantics and adapter specializations share geometric structure without joint training.

    \item \textbf{Zero-shot routing works:} IPCR achieves 95\% of oracle performance on held-out task families, validating that adapter banks can be accessed without validation examples.

    \item \textbf{Routing is lexically anchored:} The mechanism depends on task-indicative keywords rather than deep semantic invariants. Paraphrase perturbations cause 24\% routing inconsistency; keyword masking causes 44\% accuracy drops.
\end{enumerate}

This third finding is as important as the first two---it defines the boundary conditions for IPCR deployment. The method is appropriate for controlled instruction interfaces (APIs, templates) but requires robustification for open-ended queries.

\subsection{Future Directions}

Our results motivate several directions grounded in experimental evidence:

\textbf{Robust router training:} The lexical dependence revealed by H-M2 suggests paraphrase augmentation during probe training, or contrastive losses that encourage paraphrase invariance. The goal is maintaining 95\% oracle performance while improving cosine similarity under paraphrase from 0.78 to $\geq$0.90.

\textbf{Encoder alternatives:} The mean-pooling limitation of MiniLM suggests testing encoders with stronger compositional semantics (E5-large, Instructor-XL). These may provide paraphrase invariance that MiniLM lacks.

\textbf{Hybrid routing:} Combining embedding-based routing with keyword confidence could maintain high accuracy on standard instructions while providing graceful degradation on paraphrased variants. When embedding confidence is low, fall back to keyword-based routing or uniform averaging.

\textbf{Full oracle validation:} H-E1 used task names as oracle proxy. Computing per-adapter loss for each sample would provide ground-truth adapter selection labels, enabling tighter performance bounds.

The finding that zero-shot routing is achievable---but fragile---opens a research direction at the intersection of adapter composition and robust representation learning. We hope this work encourages investigation of intrinsic alignment properties in instruction-tuned models, and motivates routing architectures that preserve both efficiency and robustness.
